%!TEX TS-program = xelatex
\documentclass[12pt, a4paper, oneside]{article}

% Можно вставить разную преамбулу
\usepackage{amsmath,amsfonts,amssymb,amsthm,mathtools}  % пакеты для математики

\usepackage[english, russian]{babel} % выбор языка для документа
\usepackage[utf8]{inputenc} % задание utf8 кодировки исходного tex файла
\usepackage[X2,T2A]{fontenc}        % кодировка

% Основные шрифты 
\usepackage{fontspec}         
\setmainfont{Linux Libertine O}  % задаёт основной шрифт документа

% Математические шрифты 
\usepackage{unicode-math}     
\setmathfont[math-style=upright]{euler.otf} 

\setmathfont[range={\mathbb, \mathop, \heartsuit, \angle, \smile, \varheartsuit}]{Asana-Math.otf}

%%%%%%%%%% Работа с картинками %%%%%%%%%
\usepackage{graphicx}                  % Для вставки рисунков
\usepackage{graphics}
\graphicspath{{images/}}               % можно указать папки с картинками
\usepackage{wrapfig}                   % Обтекание рисунков и таблиц текстом

%%%%%%%%%%%%%%%%%%%%%%%% Графики и рисование %%%%%%%%%%%%%%%%%%%%%%%%%%%%%%%%%
\usepackage{tikz, pgfplots}  % язык для рисования графики из latex'a

%%%%%%%%%% Гиперссылки %%%%%%%%%%
\usepackage{xcolor}              % разные цвета

\usepackage{hyperref}
\hypersetup{
	unicode=true,           % позволяет использовать юникодные символы
	colorlinks=true,       	% true - цветные ссылки, false - ссылки в рамках
	urlcolor=blue,          % цвет ссылки на url
	linkcolor=red,          % внутренние ссылки
	citecolor=green,        % на библиографию
	pdfnewwindow=true,      % при щелчке в pdf на ссылку откроется новый pdf
	breaklinks              % если ссылка не умещается в одну строку, разбивать ли ее на две части?
}

\usepackage{todonotes} % для вставки в документ заметок о том, что осталось сделать
% \todo{Здесь надо коэффициенты исправить}
% \missingfigure{Здесь будет Последний день Помпеи}
% \listoftodos --- печатает все поставленные \todo'шки


% Опции для настройки внешнего вида 
\usepackage[paper=a4paper, top=20mm, bottom=15mm,left=20mm,right=15mm]{geometry}
\usepackage{indentfirst}       % установка отступа в первом абзаце главы

\usepackage{setspace}
\setstretch{1.15}  % Межстрочный интервал
\setlength{\parskip}{4mm}   % Расстояние между абзацами
% Разные длины в латехе https://en.wikibooks.org/wiki/LaTeX/Lengths

% мои цвета https://www.artlebedev.ru/colors/
\definecolor{titleblue}{rgb}{0.2,0.4,0.6} 
\definecolor{blue}{rgb}{0.2,0.4,0.6} 
\definecolor{red}{rgb}{1,0,0.2} 
\definecolor{green}{rgb}{0,0.6,0} 
\definecolor{purp}{rgb}{0.4,0,0.8} 

% цвета из geogebra 
\definecolor{litebrown}{rgb}{0.6,0.2,0}
\definecolor{darkbrown}{rgb}{0.75,0.75,0.75}

% меняю оформление секций 
\usepackage{titlesec}
\usepackage{sectsty}

% меняю цвет на синий
\sectionfont{\color{titleblue}}
\subsectionfont{\color{titleblue}}

% Гиперссылки
\usepackage{xcolor}   % разные цвета

\usepackage{hyperref}
\hypersetup{
	unicode=true,           % позволяет использовать юникодные символы
	colorlinks=true,       	% true - цветные ссылки
	urlcolor=blue,          % цвет ссылки на url
	linkcolor=black,          % внутренние ссылки
	citecolor=green,        % на библиографию
	breaklinks              % если ссылка не умещается в одну строку, разбивать её на две части?
}

% эпиграфы
\usepackage{epigraph}
\setlength\epigraphwidth{0.4\textwidth}
\setlength\epigraphrule{0pt}

% бульпоинты в списках
\usepackage{enumitem}
\newcommand*{\MyPoint}{\tikz \draw [baseline, fill=blue,draw=blue] circle (2.5pt);}
\renewcommand{\labelitemi}{\MyPoint}

% расстояние в списках
\setlist[itemize]{parsep=0.4em,itemsep=0em,topsep=0ex}
\setlist[enumerate]{parsep=0.4em,itemsep=0em,topsep=0ex}

% Математические операторы первой необходимости:
\DeclareMathOperator{\sgn}{sign}
\DeclareMathOperator*{\argmin}{arg\,min}
\DeclareMathOperator*{\argmax}{arg\,max}
\DeclareMathOperator{\Cov}{Cov}
\DeclareMathOperator{\Var}{Var}
\DeclareMathOperator{\Corr}{Corr}
\DeclareMathOperator{\E}{\mathop{E}}
\DeclareMathOperator{\Med}{Med}
\DeclareMathOperator{\Mod}{Mod}
\DeclareMathOperator*{\plim}{plim}

\DeclareMathOperator{\logloss}{logloss}
\DeclareMathOperator{\softmax}{softmax}

\DeclareMathOperator{\tr}{tr}

% команды пореже
\newcommand{\const}{\mathrm{const}}  % const прямым начертанием
\newcommand{\iid}{\sim i.\,i.\,d.}  % ну вы поняли...
\newcommand{\fr}[2]{\ensuremath{^{#1}/_{#2}}}   % особая дробь
\newcommand{\ind}[1]{\mathbbm{1}_{\{#1\}}} % Индикатор события
\newcommand{\dx}[1]{\,\mathrm{d}#1} % для интеграла: маленький отступ и прямая d

% одеваем шапки на частые штуки
\def \hb{\hat{\beta}}
\def \hs{\hat{s}}
\def \hy{\hat{y}}
\def \hY{\hat{Y}}
\def \he{\hat{\varepsilon}}
\def \hVar{\widehat{\Var}}
\def \hCorr{\widehat{\Corr}}
\def \hCov{\widehat{\Cov}}

% Греческие буквы
\def \a{\alpha}
\def \b{\beta}
\def \t{\tau}
\def \dt{\delta}
\def \e{\varepsilon}
\def \ga{\gamma}
\def \kp{\varkappa}
\def \la{\lambda}
\def \sg{\sigma}
\def \tt{\theta}
\def \Dt{\Delta}
\def \La{\Lambda}
\def \Sg{\Sigma}
\def \Tt{\Theta}
\def \Om{\Omega}
\def \om{\omega}

% Готика
\def \mA{\mathcal{A}}
\def \mB{\mathcal{B}}
\def \mC{\mathcal{C}}
\def \mE{\mathcal{E}}
\def \mF{\mathcal{F}}
\def \mH{\mathcal{H}}
\def \mL{\mathcal{L}}
\def \mN{\mathcal{N}}
\def \mU{\mathcal{U}}
\def \mV{\mathcal{V}}
\def \mW{\mathcal{W}}

% Жирные буквы
\def \mbb{\mathbb}
\def \RR{\mbb R}
\def \NN{\mbb N}
\def \ZZ{\mbb Z}
\def \PP{\mbb{P}}
\def \QQ{\mbb Q}


% для условия задачек 
\newcounter{problem}
\renewcommand{\theproblem}{\arabic{problem}}
\newcommand{\problemname}{\color{blue} Задача}

\newenvironment{problem}[1]{
	\addtocounter{problem}{1}\noindent{\large\bfseries \problemname{} \theproblem \, (#1 баллов).
		}
}{
	\par\bigskip
}


\usepackage{amssymb,fge}
\newcommand{\mysetminus}{\mathbin{\fgebackslash}}

\title{Yet Another Math for DS Course \\ Домашка №1}
\date{}
\author{Матричное дифФфФфФференцирование \\ (продвинутая группа)}

\begin{document}

\maketitle

\epigraph{ $\bigg(\text{ \includegraphics[scale=0.08]{tree1.png}} \bigg)^T = $  \includegraphics[scale=0.08]{tree2.png} }{<<Джек и бобовый стебель>> (1890)}

Добро пожаловать в первую домашку. Для вашего удобства она разбита на три части. В первой находятся совсем простые задачи. Если вы решите их идеально, вы наберёте $50$ баллов из $100.$ Решая только первый раздел из каждой домашки, вы будете уверенно двигаться к троечке. Если вы претендуете на большее, для вас есть разделы с более сложными задачами. 

Решение работы нужно сдать в виде pdf-файла. Решения должны быть оформлены на листочке аккуратным почерком либо затеханы на компьютере. Если у вас плохой почерк, домашка должна быть затехана. Затехать домашку можно в overleaf, typora, colab или другом любом удобном для вас сервисе.

\section*{Задачи на троечку}

Если вы идеально решаете все задачи из этого раздела, вы получаете $50$ баллов из $100$.

\begin{problem}{30}
    Найдите производные $\nabla_X f(X)$ следующих функций
    
    \begin{enumerate}
        \item  $f(X) = \tr(AX^TXBX^{-T}),$ где  $A, X, B \in \RR^{n \times n}$ 
        
        \item  $f(X) = \tr(AXB),$ где $A \in \RR^{p \times m}, B \in \RR^{n \times p}, X  \in \RR^{m \times n}.$ 
        
        \item  $f(X) = \det(X^TAX),$ где  $A, X \in \RR^{n \times n}.$ 
    \end{enumerate}
\end{problem}

\begin{problem}{10}
    Пусть $f(X) = \ln \det X,$ где $X \in \mathbb{R}^{n\times n}$. Найдите производную $\nabla_X f(X)$.
\end{problem}

\newpage

\begin{problem}{10}
Решите следующую задачу матричной оптимизации. При решении считайте, что матрица $A$ положительно определена.

\[
f(x) = x^T A x - x^Tb + c \to \min_{x}
\]
\end{problem}


\section*{Задачи на хор}

Если вы идеально решаете ещё и этот раздел, вы получаете $70$ баллов из $100$.

\begin{problem}{20}
    Предположим, что мы хотим обучить ridge-регрессию, но в качестве регуляризатора решили взять видоизменённую норму ($S$ -- некоторая фиксированная симметричная положительно-определённая матрица):
    
    \[
    L(w) = (y - Xw)^T \cdot (y-Xw) + ||Sw||^2_2.
    \]
    
    Найдите градиент $\nabla_w L(w)$ получившейся функции потерь. Чему будет равно оптимальное значение весов в этом случае? Проверьте, что найденное значение действительно является точкой минимума. 
\end{problem}


\section*{Задачи на отл}

Если вы идеально решаете ещё и этот раздел, вы получаете $90$ баллов из $100$.

\begin{problem}{10}
    Рассмотрим симметричную матрицу $A \in \mathbb{R}^{n \times n}$ и ее спектральное разложение $A = Q \Lambda Q^T$.
    Пусть $\lambda \in \mathbb{R}^n$ - это диагональ матрицы $\Lambda$ (то есть вектор, составленный из собственных значений $A$).
    Найдите производные:
    
    \begin{enumerate}
        \item $\nabla_{\lambda} \tr(A)$
        \item $\nabla_Q \tr(A)$
    \end{enumerate}
\end{problem}

\begin{problem}{10}
Пусть $x \in \RR^n \mysetminus \{0\}.$ Найдите $\nabla_x f(x)$ для \[f(x) = \langle x, x \rangle^{\langle x, x \rangle}.\]
\end{problem}


\section*{Задачи на десяточку}

Если вы идеально решаете ещё и этот раздел, вы выбиваете $100$ из $100$. Вы большой молодец. 

\begin{problem}{10}
    Рассмотрим целевую функцию логистической регрессии 
    \[
    Q(w) = \frac{1}{n} \sum_{i=1}^{n} \log (1 + \exp(-y_i \, \langle w, x_i\rangle)),
    \]
    
    \begin{enumerate} 
        \item Найдите градиент $\nabla Q_w$ и упростите итоговое выражение таким образом, чтобы в нём участвовала сигмоидная функция 
        $$\sigma(z) = \frac{1}{1 + \exp(-z)}.$$ При решении данной задачи вам может понадобиться следующий факт (убедитесь, что он действительно выполняется):
        $$\sigma'(z) = \sigma(z) (1- \sigma(z)).$$ 
        \item Выпишите, как будет выглядеть шаг градиентного спуска.
        \item Найдите вторую производную целевой функции по $w$.
        \item Выпишите квадратичную аппроксимацию для $Q(w)$ в окрестности $w=0$. Для этого разложите функцию потерь в ряд Тейлора до второго члена в окрестности точки $w=0$. С какой задачей совпадает задача минимизации квадратичной аппроксимации?
    \end{enumerate} 
\end{problem}

\end{document}
